\section{Adelic formulation and choices}\label{subsec:adelic-packet}

The model of Lemma~\ref{lem:packet-model} also identifies the measure obtained from adelic
torus Haar probability. We give the identification for arbitrary local
type, so the lattice formulation includes all finite lattice and real
diagonalization choices in the usual adelic formulation.

Let $\mathbb A_f$ be the finite adeles of $\Q$, let
$\widehat\Z=\prod_\ell\Z_\ell$, and put
$\mathbb A_\Q=\R\times\mathbb A_f$ and
$\mathbb A_K=K\otimes_\Q\mathbb A_\Q$.
Write $\widehat R=R\otimes\widehat\Z$.
Choose a rational basis $b_1,\ldots,b_n$ of $K$, and let
$\rho(\alpha)$ denote multiplication by $\alpha$ on row coefficient
vectors. It embeds the torus
\[
 \mathbf T_K=(\operatorname{Res}_{K/\Q}\mathbb G_m)/\mathbb G_m
 \quad\hbox{in}\quad G=\operatorname{PGL}_{n/\Q}.
\]
The embedding matrix $B_\sigma=(\sigma_j(b_i))_{i,j}$ satisfies
\begin{equation}\label{eq:regular-diagonalization}
 \rho(\alpha)B_\sigma
   =B_\sigma\operatorname{diag}(\sigma_1(\alpha),\ldots,\sigma_n(\alpha)).
\end{equation}
For $g_f\in G(\mathbb A_f)$, choose a lift to
$\operatorname{GL}_n(\mathbb A_f)$ and set
\[
 \widehat M(g_f)=\widehat\Z^{\,n}g_f^{-1},\qquad
 M(g_f)=K\cap\widehat M(g_f),
\]
using the chosen rational basis to identify coefficient vectors with
elements of $K$. This is a full lattice, and its completion is
$\widehat M(g_f)$: choose an integer $d\ge1$ with
$d\widehat\Z^{\,n}\subset\widehat M(g_f)
\subset d^{-1}\widehat\Z^{\,n}$, and patch in the finite quotient
$d^{-1}\Z^n/d\Z^n$. Changing the projective lift multiplies it by a
scalar finite idele. Every such idele is a rational scalar times an
element of $\widehat\Z^\times$; the latter preserves every completed
full lattice. Consequently this change only rescales $M(g_f)$ by
a rational number.

Choose $g_\infty\in G(\R)$ with
$g_\infty^{-1}\mathbf T_K(\R)^\circ g_\infty=A_n$, using the
odd-degree identification of $G(\R)$ with $\SL_n(\R)$.
For $g=(g_\infty,g_f)$, the adelic packet is the projection of
$\mathbf T_K(\Q)\backslash\mathbf T_K(\mathbb A_\Q)g$ to
$G(\Q)\backslash G(\mathbb A_\Q)/G(\widehat\Z)$, equipped with
the pushforward of torus Haar probability.

\begin{proposition}\label{prop:adelic-packet}
The adelic torus quotient is compact. Its packet and canonical
probability agree with $\mathcal P_{K,M(g_f),\sigma'}$ and
$\mu_{K,M(g_f),\sigma'}$ for an ordering $\sigma'$ of the real
embeddings. Every prescribed local homothety type occurs this way.
Changing the rational basis or the projective lifts does not change
the packet; changing the allowed real diagonalizer only reorders the
embeddings. Thus Theorem~\ref{thm:main} includes every such adelic choice.
\end{proposition}

\begin{proof}
Write $M=M(g_f)$ and $\widehat M=\widehat M(g_f)$.
First identify the ambient quotient with $X_n$. Given a finite matrix
lift $g_f$, choose a basis matrix $c\in\operatorname{GL}_n(\Q)$
for $\Q^n\cap\widehat\Z^{\,n}g_f^{-1}$. Then
$cg_f\in\operatorname{GL}_n(\widehat\Z)$, and the corresponding
real lattice is $\Z^ncg_\infty$, normalized to covolume one.
Changing this basis or either matrix lift has no effect on that
lattice. This constructs the usual quotient identification directly.
For odd $n$, the signed real $n$th root gives
$[h]\mapsto(\det h)^{-1/n}h$, identifying
$G(\R)$ with $\SL_n(\R)$ and $G(\Z)$ with $\SL_n(\Z)$.

Multiplication matrices modulo scalars give
\begin{equation}\label{eq:idelic-torus-quotient}
 \mathbf T_K(\Q)\backslash\mathbf T_K(\mathbb A_\Q)
 \simeq K^\times\backslash\mathbb A_K^\times/\mathbb A_\Q^\times.
\end{equation}
Here the pointwise lifts can also be seen by elementary linear algebra:
over any ground field $F_0$ containing $\Q$, a matrix whose projective
class lies in $\mathbf T_K$ becomes a scalar multiple of a multiplication
operator over an algebraic closure. It therefore commutes exactly with
the regular representation already over $F_0$. The centralizer of that
representation is the algebra of multiplication operators, so the matrix
belongs to $(K\otimes_\Q F_0)^\times$, with scalar ambiguity $F_0^\times$.
Finite lifts may be chosen integral at almost all primes.

A finite idele $t_f$ acts on the completed lattice by
$\widehat M\mapsto t_f^{-1}\widehat M$. Hence its translates are
exactly the prescribed local homothety type. The stabilizer of
$\widehat M$ among finite ideles is
\[
 \{z:z\widehat M=\widehat M\}
       =(\cO(M)\otimes\widehat\Z)^\times:
\]
equality holds precisely when both $z$ and $z^{-1}$ preserve
$\widehat M$. In the projective torus the stabilizer is the image
of this group, since a representative that preserves the lattice
up to a scalar idele can be rescaled to preserve it exactly.
The unit subgroup is compact open in the finite ideles and has finite
index in $\widehat R^\times$; at each prime not dividing $q$ it equals
$R_\ell^\times$. Its projective image is compact open as well.
No invertibility of $M$ over $\cO(M)$ is used.

To describe the quotient and its measure, first divide
\eqref{eq:idelic-torus-quotient} by the image of
$\widehat R^\times$. Fractional-ideal factorization and
$\mathbb A_\Q^\times
=\Q^\times\R^\times\widehat\Z^\times$
give one real fiber $E\backslash C$ for each ordinary ideal class.
These fibers are compact by the unit theorem, and the omitted
finite-unit group is compact. Thus the torus quotient is compact.
At the stabilizer of $\widehat M$, the finite fibers instead retain
the translates $\mathcal N(I)$. The same decomposition becomes
\[
 \coprod_{[I]}(\mathcal N(I)\times C)/E,
\]
with exactly the action in Lemma~\ref{lem:packet-model}.
Haar measure pushes to uniform counting measure on each finite
unit orbit and to the stated logarithmic Haar measure on $C$.
After normalization it is therefore precisely the probability in
that lemma. For $g_\infty=[B_\sigma]$,
\eqref{eq:regular-diagonalization} identifies the resulting real
lattices with its displayed map. This proves both the packet
identification and the equality of measures, including the orbit
volume weights.

Finally, every full $M$ is obtained by choosing a basis matrix for
its completed lattice and taking its inverse as $g_f$.
Any other allowed real diagonalizer has the form $[B_\sigma dP]$,
where $d$ is diagonal and $P$ is a permutation matrix, because it
must permute the joint eigenspaces. The diagonal factor translates
the real torus and permutes its sign components, preserving its
Haar probability; the permutation reorders the embeddings.
If a rational change of basis has matrix $c$, then
$B'_\sigma=cB_\sigma$ and $\rho'=c\rho c^{-1}$, while the same
packet data have representatives $g'_\infty=cg_\infty$ and
$g'_f=cg_f$. Thus the torus is conjugated and the translated packet
is left-multiplied by $c$, which disappears in the ambient rational quotient.
\end{proof}
